\sscp{Other classificatinos}{07-01-03}
\citet{Dyen1978b} regrouped Stresemann’s Sub-Ambon closer with
Sub-Seram and farther from Sub-Buru based mostly on lexicostatistics.
Both \citeauthor{Collins1983}’ (\citeyear{Collins1983}) Nunusaku
and our \tsc{Ambon-Seram} subgroup support the conclusion of
Dyen’s adjustment, although unlike Collins we do not link \tsc{Ambon-Seram} languages
with \tsc{Sula-Buru} at a higher level.
\citet[45]{Dyen1965} would “perhaps” include Soboyo (\tsc{Taliabu})
with his Moluccan linkage.

\citet{Blust1981}, focusing on Soboyo (Taliabu Island),
grouped the languages of the Sula archipelago in the northwest with Buru.
After careful examination, the strongest basis for such a link
is that they share *R > \it{h}.
Elsewhere \citet[279]{Blust1993} suggests that Geser
and Watubela in far-eastern Seram
and offshore islands should be a “primary branch” of CMP.

In the same work \citet[276ff]{Blust1993} defined a “Yamdena-North
Bomberai” subgroup based on a bundle of lexical innovations,
most of which are not exclusive to the group of languages on
which Blust was focused.\footnote{
For example,
reflexes of *bat-bata ‘female’ are found in some \tsc{Seram-Tanimbar-Bomberai},
\tsc{Sula-Buru}, and \tsc{Timor-Babar} languages; *ndundum ‘thunder’
is found in some \tsc{Seram-Tanimbar-Bomberai},
\tsc{Aru}, \tsc{Timor-Babar}, and \tsc{Bima-Lembata} languages;
*ŋarak (which we reconstruct as **ŋaRək) ‘year’ is found in some
\tsc{Sula-Buru}, \tsc{Ambon-Seram}, \tsc{Aru}, and \tsc{Seram-Tanimbar-Bomberai}
languages; and so forth.}
Yet words like **kəRi ‘tongue’,
while found beyond Blust’s “Yamdena-North Bomberai” languages,
nevertheless do seem to be restricted to the wider STB  subgroup.
While the details of Blust’s subgroup itself is problematic,
linking Yamdena to the Austronesian languages on the Bomberai
peninsula remains an important contribution.

\citet{Mills1991} argues for a “Tanimbar-Kei” subgroup that includes
Yamdena, Fordata, and Kei.
He further linked these languages with Selaru
and Leti within a “Southeast Maluku” subgroup,
but presented very little data to support it.
His appeal to “binding” (1991:250ff) to further link southeast
Maluku to Selaru, the languages of southwest Maluku,
Timor and west Sumba are not persuasive
(see \citeauthor{Blust1993} [\citeyear[276]{Blust1993}], and \srf{13-02-03}).
\citet{Mills1991} did not investigate the possibility of linking
Kei, Fordata and Yamdena with languages to the north.

\citet{GrayGreenhill2009} include the languages of our \tsc{Taliabu},
\tsc{Sula-Buru}, \tsc{Ambon-Seram}, \tsc{Aru} subgroups, and
part of our \tsc{Seram-Tanimbar-Bomberai} subgroup under a single
subgroup they call “Central Maluku” (p.480).
Curiously their algorithm places Buru
and Soboyo as even more divergent from the other languages of
their Central Maluku than are the \tsc{Aru} languages (\crf{ch:Aru}),
whereas \citet{Collins1981,Collins1983} repeatedly claims Buru is “closely
related” to Ambon and Seram languages, and by association,
so is Soboyo and Taliabu
(\citeyear[15, 19--20]{Collins1983}, \citeyear[86, footnote 1]{Collins1989}).
These two views are not compatible
and raise many questions.
We believe our classification based primarily on the historical
phonology provides a solution. \citet{GrayGreenhill2009}
also group their “Yamdena-North Bomberai” separately from
their Central Maluku subgroup.
We place those languages within our \tsc{Seram-Tanimbar-Bomberai} subgroup.

\citet{Chlenov1976}, working mostly with lexicostatistics,
included the languages of East Seram,
but not the languages of the Sula archipelago,
within his “Ambon” (roughly Collins’ Central Maluku) subgroup.
He also placed Banda as its own node within his Ambon subgroup.
Diverging from several others,
Chlenov placed Geser and Watubela as their own node under his
separate “South Moluccan” subgroup,
linking them southward to languages such as Kei, Fordata and Yamdena.
This proposal is in alignment with our \tsc{Seram-Tanimbar-Bomberai}.

\citet{Taguchi1989} tabulates the lexical similarity for languages
in west Seram, central Seram, parts of east Seram,
and in the west,
several languages in the Manipa Strait.
His finding of Boano as being a bit divergent from the rest parallels
our findings in the historical phonology.

\citet{LoskiLoski1989} also working mostly with lexicostatistics,
tabulate the lexical similarity for languages from central Seram,
east Seram, and offshore islands including Geser-Gorom,
Watubela, \tsc{Teor-Kur}, and Kei.
Their findings \citeyear[112]{LoskiLoski1989} of Banda (Elat) being equally
unrelated to everything else mirror our placement of Banda,
based primarily on the historical phonology,
as an isolate in the region.\footnote{\label{fn:BandaLexicalSimilarity}
		While \citet{Collins1986}
		groups Banda together with Geser
		and Watubela under his Proto-Banda,
		\citet[112]{LoskiLoski1989} find Banda -- Geser-Gorom =
		39{\%}, Banda -- Watubela = 39{\%},
		Banda with \tsc{East Seram} languages (Masiwang,
		Bobot) = 35--36{\%}, Banda -- \tsc{Teor-Kur} = 40{\%},
		Banda -- Kei = 38{\%}.}
(See \crf{ch:Ban}.)

Focusing further south, \citet{Hughes1987} lexicostatistics study
was restricted geographically in scope to Kei, Tanimbar and Aru.
He did not investigate how some of those languages might be related
to the languages north of Kei,
extending to eastern Seram,
a very remote region from any direction.